\documentclass{article}
\usepackage[utf8]{inputenc}
\usepackage{amsmath}
\usepackage{amssymb}
\usepackage{url}
\usepackage{hyperref}

\title{Derivación de la ecuación de Tsiolkovsky para un cohete}
\author{Franz Chouly}
\date{\today}

\begin{document}

\maketitle

\section{Introducción}
La ecuación
\[
m(t) \frac{dv}{dt} = -u \frac{dm}{dt} - m(t) g,
\]
describe el movimiento de un cohete que expulsa masa (combustible) a alta velocidad. Es conocida como la \textbf{ecuación de movimiento para sistemas de masa variable} y fue desarrollada por \textbf{Konstantin Tsiolkovsky} (1857--1935), pionero de la astronautica teórica.

\section{Origen histórico}
\begin{itemize}
    \item \textbf{Tsiolkovsky (1903)}: Publicó su obra ``Investigación de los Espacios Cósmicos por Medio de Aparatos de Reacción'', donde derivó la ecuación por primera vez.
    \item \textbf{Robert Goddard (1920s)}: Implementó experimentalmente los principios de Tsiolkovsky con los primeros cohetes de combustible líquido.
\end{itemize}

\section{Derivación}
\subsection{Hipótesis}
\begin{enumerate}
    \item Velocidad de expulsión $u$ constante (relativa al cohete).
    \item Masa del cohete $m(t)$ disminuye: $\frac{dm}{dt} < 0$.
    \item Fuerza externa: gravedad $-m(t)g$ (en dirección opuesta al movimiento).
\end{enumerate}

\subsection{Paso a Paso}

1. \textbf{Momento lineal total} en $t + dt$:
   \[
   p(t + dt) = (m + dm)(v + dv) + (-dm)(v - u).
   \] 
Simplificando (ignorando $dm \cdot dv$):
   \[
   p(t + dt) = mv + m \, dv + u \, dm.
   \]
2. \textbf{Variación del momento}:
   \[
   dp = p(t + dt) - p(t) = m \, dv + u \, dm.
   \]
   Dividiendo por $dt$:
   \[
   \frac{dp}{dt} = m \frac{dv}{dt} + u \frac{dm}{dt}.
   \]
3. \textbf{Segunda Ley de Newton}: la fuerza externa (gravedad) es igual a $\frac{dp}{dt}$.
   \[
   F_{\text{ext}} = -m(t) g = m \frac{dv}{dt} + u \frac{dm}{dt}.
   \]
   Reordenando:
   \[
   m(t) \frac{dv}{dt} = -u \frac{dm}{dt} - m(t) g.
   \]

\subsection{Interpretación}
\begin{itemize}
    \item $-u \frac{dm}{dt}$: \textbf{Empuje} (positivo, ya que $\frac{dm}{dt} < 0$).
    \item $-m(t)g$: Fuerza gravitatoria.
\end{itemize}

\section{Ecuación de Tsiolkovsky (sin gravedad)}
Si $g = 0$ y $u$ es constante, integrando:
\[
\Delta v = u \ln \left( \frac{m_0}{m_f} \right),
\]
donde:
\begin{itemize}
    \item $\Delta v$: Cambio de velocidad (``delta-v''),
    \item $m_0$: Masa inicial (combustible + estructura),
    \item $m_f$: Masa final (solo estructura).
\end{itemize}

%\section{Conclusión}
%La ecuación surge de aplicar las leyes de Newton a sistemas de masa variable. Es fundamental para calcular trayectorias espaciales y el diseño de cohetes.

\section*{Referencias}
\begin{description}
    \item[NASA.] \textit{Beginner’s Guide to Aeronautics: Ideal Rocket Equation}. \\
    Disponible en: \url{https://www1.grc.nasa.gov/beginners-guide-to-aeronautics/ideal-rocket-equation/}. \\
    Guía introductoria sobre la ecuación de Tsiolkovsky, con explicaciones accesibles y ejemplos prácticos.

    \item[Richard, T.] (2025). \textit{SML\_Tom\_1.pdf: Notes on Mathematical Foundations for Rocket Propulsion}. \\
    Universidad de Nagoya. \\
    Disponible en: \url{https://www.math.nagoya-u.ac.jp/~richard/teaching/f2025/SML_Tom_1.pdf}. \\
    Apuntes académicos sobre fundamentos matemáticos en propulsión de cohetes.

    \item[ArXiv.] (2024). \textit{Advances in Rocket Propulsion Theory}. \\
    Disponible en: \url{https://arxiv.org/pdf/2406.00055}. \\
    Artículo reciente sobre avances teóricos en propulsión.

    \item[Tsiolkovsky, K. E.] \textit{Obras Científicas y Legado}. \\
    Instituto Tsiolkovsky. \\
    Disponible en: \url{https://www.tsiolkovsky.org/ru/nauchnoe-nasledie/}. \\
    Recopilación de los trabajos originales de Konstantin Tsiolkovsky.

    \item[Vanier College.] \textit{Rockets: A Mathematical Introduction}. \\
    Proyecto ITI. \\
    Disponible en: \url{https://gauss.vaniercollege.qc.ca/~iti/proj/Rockets.pdf}. \\
    Documento educativo que combina teoría matemática y aplicaciones prácticas.
\end{description}

%REFS
%  https://www1.grc.nasa.gov/beginners-guide-to-aeronautics/ideal-rocket-equation/

% https://www.math.nagoya-u.ac.jp/~richard/teaching/f2025/SML_Tom_1.pdf

%https://arxiv.org/pdf/2406.00055

%https://www.tsiolkovsky.org/ru/nauchnoe-nasledie/

% https://gauss.vaniercollege.qc.ca/~iti/proj/Rockets.pdf

\end{document}